\documentclass[10pt,a4paper]{amsart}
\usepackage[margin=19mm]{geometry}
\usepackage{amsmath,amssymb,mathtools}
\usepackage[scr=rsfs,cal=euler]{mathalfa}
\usepackage{xfrac}
\usepackage[unicode,hidelinks,bookmarks=false]{hyperref}
\newcommand{\ie}{i.e.\ }
\newcommand{\cf}{cf.\ }
\newcommand{\resp}{respectively\ }
\newcommand{\Subsection}{\subsection}
\providecommand{\nobreakdash}{\nobreak\hbox{-}}
\setlength{\emergencystretch}{2em}
\multlinegap0pt
\usepackage{amssymb}
\usepackage{mathtools}
\usepackage[shortlabels]{enumitem}
\usepackage{tikz}
\newtheorem{theorem}{Theorem}
\newcommand{\R}{\mathbb{R}}
\newcommand{\e}{\mathrm{e}}
\newcommand{\supp}{\operatorname{supp}}
\title{Sharp $L^2$ Estimates for $(2+1)$-dimensional oscillatory integral operators with homogeneous binomial phases}
\author{Chu-hee Cho, Jin Bong Lee,, Chan Woo Yang}
\date{Source: arXiv:2608.08683v1 (CC BY 4.0)}
\begin{document}
\maketitle

\noindent\emph{Source and licence.} Sharp $L^2$ Estimates for $(2+1)$-dimensional oscillatory integral operators with homogeneous binomial phases. Version arXiv:2608.08683v1 (CC BY 4.0), distributed by arXiv under the Creative Commons Attribution 4.0 International licence, http://creativecommons.org/licenses/by/4.0/, as recorded in the arXiv licence metadata for this exact version and shown on the abstract page https://arxiv.org/abs/2608.08683v1. Full licence notice: \texttt{licenses/arxiv-2608.08683-CC-BY-4.0.txt}.\par\smallskip

\noindent\textit{Editorial scope.} Counted unit: the article's theorem thm:lbound (source0.tex lines 342--353). The article's complete original proof of it is delivered with it (source0.tex lines 867--1110, one \texttt{proof} environment of 244 lines, which the article places away from the statement). No mathematics is written, repaired or re-derived: the statement and the proof are the article's own lines, in the article's own notation, and no retained line is substituted or re-flowed.  Retained uncounted as the source-local context the proof needs: lines 200--204; lines 217--222; lines 238--240.  Nothing was omitted from any retained span, so the package carries no omission record.  No retained line contains a citation, so the wrapper carries no bibliography and no bibliography entry.  No retained line contains a figure environment, an image-inclusion command or drawing code, so no image is delivered and the package declares no asset dependency.  Editorial formatting only, in addition: an AMS \texttt{amsart} preamble that restates the article's own declarations and macros used by the retained lines, namely the environments \texttt{theorem} on separate counters, together with the standard packages those lines need.  The retained environments print in this document as Theorem 1 in order of appearance, whereas the article prints the same items with the numbering of its own full document.  The complete native source of the article as distributed in the arXiv e-print for this exact version is bundled unchanged as \texttt{sources/207181/source0.tex}.
    \begin{equation}\label{eq:intro_phase}
    \Phi(x,y,t)=x^{k-k_P}t^{k_P}+y^{k-k_Q}t^{k_Q},
    \qquad
    1\le k_P<k_Q<k.
    \end{equation}

    \begin{equation}\label{eq:intro_operator}
    T_\lambda f(x,y)=
    \int_{\mathbb R}
    e^{i\lambda \Phi(x,y,t)}
    a(x,y,t)f(t)\,\mathrm dt,
    \end{equation}

		Let \(k,k_P,k_Q\) satisfy \begin{equation}\label{eq:intro_parameters}
		k\ge 3,\qquad 1\le k_P<k_Q<k.	
	\end{equation} Define

	\begin{theorem}\label{thm:lbound}
		Let \(k,k_P,k_Q\) satisfy \eqref{eq:intro_parameters}, and let \(T_\lambda\) be defined by
		\eqref{eq:intro_operator} with phase \eqref{eq:intro_phase}.  Assume in addition that
		\begin{equation*}
			a(0,0,0)\neq 0.
		\end{equation*}
		Then there exist constants \(c>0\) and \(\lambda_1\ge2\) such that, for all \(\lambda\ge\lambda_1\),
		\begin{equation}\label{eq:lowerbound_main}
			\|T_\lambda\|_{2\to2}
			\ge c\,\lambda^{-\rho_*}.
		\end{equation}
	\end{theorem}

	\begin{proof}[Proof of Theorem~\ref{thm:lbound}]
	    \iffalse Let $\theta_0\in[0,2\pi)$ be chosen so that
		\[
		e^{-i\theta_0}a(0,0,0)=|a(0,0,0)|>0.
		\]
		Replacing $a$ by $e^{-i\theta_0}a$ does not change $\|T_\lambda\|_{2\to2}$.
        \fi
		Since the estimate is invariant under replacing $a$ by $-a$,
        we may assume
        \[
        a(0,0,0)>0.
        \]
		By continuity, there exist constants $r_0\in(0,1]$ and $c_0\in(0,1]$ such that
		$a(x,y,t)\neq0$ on $\{|x|,|y|,|t|\le r_0\}$ and
		\begin{equation}\label{eq:amp_sector}
			|a(x,y,t)|\ge c_0
			\qquad\text{whenever }|x|,|y|,|t|\le r_0.
		\end{equation}
		
		Fix a constant $r_1\in(0,r_0]$ to be specified below and let $\delta\ge 0$.
		Define
		\begin{equation*}
			f_\lambda(t):=\mathbf{1}_{[-r_1\lambda^{-\delta},\,r_1\lambda^{-\delta}]}(t).
		\end{equation*}
		Then
		\begin{equation}\label{eq:test_L2_norm}
			\|f_\lambda\|_{L^2(\R)}=(2r_1)^{1/2}\lambda^{-\delta/2}.
		\end{equation}
		Let $c\in(0,r_0]$ be a small constant to be chosen below and define, for $\lambda\ge 2$,
		\begin{equation}\label{eq:Omega_def}
			\Omega_\lambda(\delta)
			:=\Bigl\{(x,y)\in\R^2:\ |x|\le c\,X_\lambda(\delta),\ |y|\le c\,Y_\lambda(\delta)\Bigr\},
		\end{equation}
		where
		\begin{equation}\label{eq:XY_def}
			X_\lambda(\delta):=\min\{\lambda^{(k_P\delta-1)/(k-k_P)},\,1\},\qquad
			Y_\lambda(\delta):=\min\{\lambda^{(k_Q\delta-1)/(k-k_Q)},\,1\}.
		\end{equation}
		By construction, $|x|,|y|\le c\le r_0$ for all $(x,y)\in\Omega_\lambda(\delta)$.
		Moreover, since $r_1\le r_0$ and $\delta\ge0$, we also have $|t|\le r_0$ for all
		$t\in\supp f_\lambda$ and all $\lambda\ge2$.
		Therefore, for $(x,y)\in\Omega_\lambda(\delta)$ and $t\in\supp f_\lambda$ we may invoke
		\eqref{eq:amp_sector}.
		The key point is the following estimate:
		\begin{align}\label{T:lbound}
			|T_\lambda f_\lambda(x,y)| \gtrsim \lambda^{-\delta} \mathbf{1}_{\Omega_\lambda(\delta)}(x,y).
		\end{align}
		

		Indeed, for $(x,y)\in\Omega_\lambda(\delta)$ and $|t|\le r_1\lambda^{-\delta}$, we have
		\begin{align*}
			\lambda|x^{k-k_P}t^{k_P}|
			&\le \lambda\,(cX_\lambda(\delta))^{k-k_P}\,(r_1\lambda^{-\delta})^{k_P}.
		\end{align*}
		If $\delta\le 1/k_P$, then
		\[
		X_\lambda(\delta)=\lambda^{(k_P\delta-1)/(k-k_P)},
		\]
		hence
		\[
		\lambda\,(cX_\lambda(\delta))^{k-k_P}\,(r_1\lambda^{-\delta})^{k_P}
		=\lambda\,c^{k-k_P}\lambda^{k_P\delta-1}\,r_1^{k_P}\lambda^{-k_P\delta}
		=c^{k-k_P}r_1^{k_P}.
		\]
		If $\delta\ge 1/k_P$, then $X_\lambda(\delta)=1$, hence
		\[
		\lambda\,(cX_\lambda(\delta))^{k-k_P}\,(r_1\lambda^{-\delta})^{k_P}
		\le \lambda\,c^{k-k_P}\,r_1^{k_P}\lambda^{-k_P\delta}
		=c^{k-k_P}r_1^{k_P}\lambda^{1-k_P\delta}
		\le c^{k-k_P}r_1^{k_P}.
		\]
		Thus, in all cases,
		\begin{equation}\label{eq:phase_term1}
			\lambda|x^{k-k_P}t^{k_P}|\le c^{k-k_P}r_1^{k_P}.
		\end{equation}
		By the same argument, using $Y_\lambda(\delta)$ and $k_Q$ in place of $X_\lambda(\delta)$ and $k_P$,
		\begin{equation}\label{eq:phase_term2}
			\lambda|y^{k-k_Q}t^{k_Q}|\le c^{k-k_Q}r_1^{k_Q}.
		\end{equation}
		Combining \eqref{eq:phase_term1}--\eqref{eq:phase_term2} yields
		\[
		|\lambda\Phi(x,y,t)|
		\le c^{k-k_P}r_1^{k_P}+c^{k-k_Q}r_1^{k_Q}.
		\]
		Choose $r_1\in(0,r_0]$ and then $c\in(0,r_0]$ so that
		\[
		c^{k-k_P}r_1^{k_P}+c^{k-k_Q}r_1^{k_Q}\le \frac{\pi}{20}.
		\]
		This proves \begin{equation*}
		 	|\lambda\Phi(x,y,t)|\le \frac{\pi}{20}
		 	\qquad\text{whenever }(x,y)\in\Omega_\lambda(\delta)\text{ and }t\in\supp f_\lambda.
		 \end{equation*}
		
	
		Together with \eqref{eq:amp_sector} we obtain
		\begin{equation*}
			\e^{\,i\lambda\Phi(x,y,t)}a(x,y,t)\bigr)
			\gtrsim c_0
			\qquad \text{for all }t\in\supp f_\lambda.
		\end{equation*}
		\iffalse Hence for every such $t$,
		\[
		\Re\bigl(\e^{\,i\lambda\Phi(x,y,t)}a(x,y,t)\bigr)
		\ge
		\cos\Bigl(\frac{1}{10}\Bigr)\,
		\bigl|\e^{\,i\lambda\Phi(x,y,t)}a(x,y,t)\bigr|
		=
		\cos\Bigl(\frac{1}{10}\Bigr)\,|a(x,y,t)|.
		\]
        \fi
		Since
        $f_\lambda\ge0$
        and
        $\Re(e^{i\lambda\Phi}a)\gtrsim c_0$
        on the support,, we conclude that for all
		$(x,y)\in\Omega_\lambda(\delta)$,
		\begin{align}
			|T_\lambda f_\lambda(x,y)|
			&\ge
			 \int_{\supp (f_\lambda)}|a(x,y,t)|\,\mathrm dt\nonumber\\
			&\gtrsim c_0\,|\supp (f_\lambda)|\nonumber\\
			&= c_0\,(2r_1)\lambda^{-\delta}
			\ \gtrsim\ \lambda^{-\delta}.\label{eq:pointwise_lower}
		\end{align}
		By \eqref{eq:pointwise_lower}, we obtain \eqref{T:lbound}.
		
		Integrating \eqref{T:lbound} over $x,y$ yields
		\begin{equation}\label{eq:L2_lower_raw}
			\|T_\lambda f_\lambda\|_{L^2(\R^2)}
			\ge \|T_\lambda f_\lambda\|_{L^2(\Omega_\lambda(\delta))}
			\gtrsim \lambda^{-\delta}\,|\Omega_\lambda(\delta)|^{1/2}.
		\end{equation}
		Combining \eqref{eq:L2_lower_raw} with \eqref{eq:test_L2_norm} gives
		\begin{equation}\label{eq:op_lower_ratio}
			\|T_\lambda\|_{2\to2}
			\ge \frac{\|T_\lambda f_\lambda\|_2}{\|f_\lambda\|_2}
			\gtrsim \lambda^{-\delta/2}\,|\Omega_\lambda(\delta)|^{1/2}.
		\end{equation}
		By \eqref{eq:Omega_def}--\eqref{eq:XY_def}, we have
		\begin{equation}\label{eq:Omega_measure}
			|\Omega_\lambda(\delta)| \approx X_\lambda(\delta)\,Y_\lambda(\delta),
		\end{equation}
		with constants depending only on $c$.
		Thus, the lower bound in \eqref{eq:op_lower_ratio} is determined by the size of
		$X_\lambda(\delta)$ and $Y_\lambda(\delta)$, leading to three regimes.
		
		\smallskip
		\noindent\textbf{Case 1: $0\le \delta\le 1/k_Q$.}
		Then
		\[
		X_\lambda(\delta)=\lambda^{(k_P\delta-1)/(k-k_P)},
		\qquad
		Y_\lambda(\delta)=\lambda^{(k_Q\delta-1)/(k-k_Q)}.
		\]
		By \eqref{eq:op_lower_ratio}--\eqref{eq:Omega_measure},
		\begin{align}
			\|T_\lambda\|_{2\to2}
			&\gtrsim \lambda^{-\delta/2}\,
			\lambda^{\frac12\bigl(\frac{k_P\delta-1}{k-k_P}+\frac{k_Q\delta-1}{k-k_Q}\bigr)}\nonumber\\
			&=\lambda^{-\rho_1-\frac{\delta}{2}\bigl(1-\frac{k_P}{k-k_P}-\frac{k_Q}{k-k_Q}\bigr)}.\label{eq:case1_bound}
		\end{align}
		The exponent in \eqref{eq:case1_bound} is affine in $\delta$. Hence the optimum is attained at one of the endpoints.
		If
		\[
		1-\frac{k_P}{k-k_P}-\frac{k_Q}{k-k_Q}\ge 0,
		\]
		then the lower bound in \eqref{eq:case1_bound} is optimized at $\delta=0$, giving
		\[
		\|T_\lambda\|_{2\to2}\gtrsim \lambda^{-\rho_1}.
		\]
		If
		\[
		1-\frac{k_P}{k-k_P}-\frac{k_Q}{k-k_Q}<0,
		\]
		then the lower bound in \eqref{eq:case1_bound} is optimized at $\delta=1/k_Q$, giving
		\[
		\|T_\lambda\|_{2\to2}\gtrsim \lambda^{-\rho_2}.
		\]
		
		\smallskip
		\noindent\textbf{Case 2: $1/k_Q\le \delta\le 1/k_P$.}
		Then $Y_\lambda(\delta)=1$ while
		\[
		X_\lambda(\delta)=\lambda^{(k_P\delta-1)/(k-k_P)}.
		\]
		Hence
		\begin{align}
			\|T_\lambda\|_{2\to2}
			&\gtrsim \lambda^{-\delta/2}\,\lambda^{\frac12\cdot\frac{k_P\delta-1}{k-k_P}}
			=\lambda^{-\frac1{2(k-k_P)}-\frac{\delta}{2}\bigl(1-\frac{k_P}{k-k_P}\bigr)}.\label{eq:case2_bound}
		\end{align}
		Since
		\[
		1-\frac{k_P}{k-k_P}=\frac{k-2k_P}{k-k_P},
		\]
		the exponent in \eqref{eq:case2_bound} is affine in $\delta$. Again,
        the exponent is affine in $\delta$,
        so it suffices to compare the two endpoints:
		\begin{itemize}
			\item if $k>2k_P$, the lower bound is optimized at $\delta=1/k_Q$, yielding
			\[
			\|T_\lambda\|_{2\to2}\gtrsim \lambda^{-\rho_2};
			\]
			\item if $k\le 2k_P$, the lower bound is optimized at $\delta=1/k_P$, yielding
			\[
			\|T_\lambda\|_{2\to2}\gtrsim \lambda^{-\rho_3}.
			\]
		\end{itemize}
		
		\smallskip
		\noindent\textbf{Case 3: $\delta\ge 1/k_P$.}
		Then $X_\lambda(\delta)=Y_\lambda(\delta)=1$.
			Taking the endpoint choice $\delta=1/k_P$ in \eqref{eq:op_lower_ratio}, we obtain
			\[
			\|T_\lambda\|_{2\to2}
			\gtrsim
			\lambda^{-1/(2k_P)}
			=
			\lambda^{-\rho_3}.
			\]
			For $\delta>1/k_P$, the lower bound $\lambda^{-\delta/2}$ is weaker, so the optimal choice in this regime is
			$\delta=1/k_P$.
		
		
	
			Collecting the three regimes,
we obtain the lower bounds at
			\[
			\delta=0,\qquad \delta=\frac1{k_Q},\qquad \delta=\frac1{k_P}
			\]
			are admissible choices in the three regimes above and produce, respectively, the lower bounds
			\[
			\lambda^{-\rho_1},\qquad
			\lambda^{-\rho_2},\qquad
			\lambda^{-\rho_3}.
			\]
			Taking the strongest among them yields, for all sufficiently large $\lambda$,
		\[
		\|T_\lambda\|_{2\to2}
		\gtrsim \lambda^{-\min\{\rho_1,\rho_2,\rho_3\}}
		=\lambda^{-\rho_*},
		\]
		which is \eqref{eq:lowerbound_main}.
	\end{proof}

\end{document}
